\scanpage{26}
% L01–06
\uline{Proposition} Soient $u,u'$ deux sections de $G$, \deleted{telles} \inserted{on} suppose $A(u)-2$ \inserted{et $A(u)+2$} inv. [i.e. $u_s^2$ non unipotent $\forall s\in S$]. Pour que $u$ et \unsure{$v$|$u'$} soient conjugués loc. fpqc (ou encore, conjugués si $G$ forme de $\mathrm{GP}(1)$) il f. et \uuline{s.} que $A(u)=A(u')$.
\margin{\unsure{Ceci|Celà} reste vrai sans restriction sur \unsure{$\alpha(u)$|$A(u)$}, si $S$ spectre d'un corps et $u_s$ non central \unsure{$K_s$|$\forall s$})}

% L07–12
\deleted{En effet,} \inserted{Néc. évidente, prouvons suff.} L'hypothèse \inserted{sur $u$ et $u'$} implique que $u,v$ sont contenus chacun dans un (unique) tore maximal. Utilisant le th. de conjugaison, on peut supposer que c'est le $=$, et par localisation, que c'est $G_m$. Donc $u,u'$ correspondent à des sections $\lambda,\lambda'$ de $G_m$, et l'hypo $A(u)=A(u')$ signifie

% L13–18
\[
\lambda+\lambda^{-1}=\lambda'+\lambda'^{-1}\qquad\text{soit }\alpha\in\Gamma(S,\unsure{\underline O_S|\mathcal O_S})\tag*{$(*)$}
\]
ce qui signifie que $\lambda$ et $\lambda'$ sont \deleted{\unsure{racines}} \inserted{solutions} de la $=$ équation
\[
\deleted{\unsure{x}}t^2-\alpha t+1=0\tag*{$(**)$}
\]
dont le discriminant est
\[
\alpha^2-4=(\underbrace{\alpha-2}_{\text{inv.}})(\alpha+2)
\]

% L19–28
\deleted{il n'est donc nul qu'aux pts où $\alpha(s)=-2$.} \inserted{donc inv. par hypo} \deleted{Si $\alpha^2-4$ inv. i.e. $\alpha+2$ inv. i.e. $\alpha(s)\ne-2$ en tt $s\in S$, alors} Cette équation définit un rev. \uline{étale} de $S$, et \deleted{on trouve que l'on a loc deux solutions distinctes} deux sections sont soit égales loc, soit conjuguées par l'involution canonique $t\longmapsto t^{-1}$, on gagne donc. Il reste à faire une étude locale \deleted{\unsure{aux}} plus serrée aux pts où $\alpha(s)=\pm2$, i.e. $\lambda_s=\lambda'_s=\pm1$, prenons le cas $-1$ p.ex.
% archive sheet number: 25
